\paragraph{Theorem T8: Graceful Degradation Dominance with Useful Work.}
Fix a finite horizon and a paired admissible fault trace. Let \(P_G\) be an
ORIUS graceful-degradation policy and \(P_U\) an uncontrolled persistence policy
evaluated on that same trace. Let \(Z_t^G,Z_t^U\in\{0,1\}\) be true-state
violation indicators, and let \(W_G,W_U\ge0\) be the corresponding useful-work
totals on the horizon. For a declared operating threshold
\(\lambda\in(0,1]\), if
\[
  Z_t^G\le Z_t^U\quad \forall t
  \qquad\text{and}\qquad
  W_G\ge \lambda W_U,
\]
then \(P_G\) weakly dominates \(P_U\) in the safety--work preorder
\[
  P\preceq_{\lambda} Q
  \quad\Longleftrightarrow\quad
  \sum_t Z_t^P\le\sum_t Z_t^Q
  \;\text{ and }\;
  W_P\ge\lambda W_Q .
\]
This is a paired-trace dominance statement, not a claim that \(P_G\) is globally
optimal among all fallback policies.

\paragraph{Proof.}
Summing the pointwise inequality over the fixed horizon gives
\[
  \sum_t Z_t^G\le \sum_t Z_t^U,
\]
so the graceful policy has no larger violation count, and therefore no larger
violation rate, on the paired trace. The second hypothesis is exactly the
useful-work lower bound \(W_G\ge\lambda W_U\). These two inequalities are the
two coordinates of the stated preorder. Immediate shutdown can satisfy the
safety coordinate, but it is a T8 witness only if it also satisfies the same
useful-work threshold; otherwise it is outside the theorem's dominance claim.
